% =============================================================================
%  Tài liệu 05 - Báo cáo theo dõi và kiểm soát (Monitoring & Control Report)
%  Build: latexmk 05-monitoring-and-control-report.tex   (pdfLaTeX, output in build/)
% =============================================================================
\documentclass[12pt,a4paper]{report}
\usepackage{sumtalk-report}
% =============================================================================
%  sumtalk-info.tex - information shared by every document (01-06).
%  Edit here once; every cover and footer picks it up.
% =============================================================================
\newcommand{\ProjectName}{SumTalk}
\newcommand{\CourseName}{QUẢN LÝ DỰ ÁN PHẦN MỀM}          % cover (upper case)
\newcommand{\CourseNameNormal}{Quản lý Dự án Phần mềm}     % footer
\newcommand{\AcademicYear}{2026 -- 2027}
\newcommand{\Instructor}{\canDien[họ tên giảng viên]}

% One \item per member: "Full name -- student ID", sorted by given name.
\newcommand{\TeamMembers}{%
  \item Nguyễn Anh Duy -- \canDien[MSSV]
  \item Đỗ Thành Đạt -- \canDien[MSSV]
  \item Trần Hải Đăng -- \canDien[MSSV]
  \item Châu Hùng Hào -- \canDien[MSSV]
  \item Bùi Đoàn Minh Nhật -- \canDien[MSSV]
  \item Võ Hoàng Sơn -- \canDien[MSSV]
}


\newcommand{\DocTitle}{BÁO CÁO THEO DÕI VÀ KIỂM SOÁT}
\newcommand{\DocTitleNormal}{Báo cáo theo dõi và kiểm soát}
\newcommand{\DocTitleEn}{Monitoring \& Control Report}
\newcommand{\DocRunningTitle}{Báo cáo theo dõi và kiểm soát}
\newcommand{\DocCode}{SUMTALK-MCR}
\newcommand{\DocVersion}{2.1 -- 24/11/2026}
\newcommand{\DocStatus}{Chờ duyệt}
\newcommand{\DocAuthor}{\canDien[Project Manager]}
\newcommand{\DocReviewers}{Project Manager, Product Owner (Sponsor nhận bản báo cáo mỗi sprint)}
\newcommand{\DocReaders}{Project Manager, Product Owner, Sponsor}
\newcommand{\CoverDate}{TP. HỒ CHÍ MINH, THÁNG 11/2026}
\hypersetup{pdftitle={SumTalk -- Báo cáo theo dõi và kiểm soát},pdfauthor={Nhóm SumTalk}}

\begin{document}

\makereportcover

% ============================================================ Front matter
\chapter*{Thông tin tài liệu}
\addcontentsline{toc}{chapter}{Thông tin tài liệu}

\docinfotable

\section*{Lịch sử phiên bản}

\begin{table}[H]
  \centering
  \begin{tabularx}{\linewidth}{@{}Q{1.6cm}Q{2.3cm}LP{3.2cm}@{}}
    \toprule
    \thead{Phiên bản} & \thead{Ngày} & \thead{Nội dung thay đổi} & \thead{Người thực hiện} \\
    \midrule
    1.0 & 12/11/2026 & Bản đầu, gồm bốn tài liệu riêng: báo cáo trạng thái sprint, Earned Value, biên bản họp và
      Retrospective, Change Log & \canDien \\
    2.0 & 13/11/2026 & Gộp thành một tài liệu, đánh lại số mục (CR-08) & \canDien \\
    2.1 & 24/11/2026 & Chuyển sang định dạng báo cáo PDF (LaTeX), sắp xếp lại thành sáu chương; vẽ biểu đồ
      burndown, đường S và quy trình kiểm soát thay đổi & \canDien \\
    \bottomrule
  \end{tabularx}
\end{table}

\needspace{8\baselineskip}
\heading{Quy ước trong tài liệu}
\begin{itemize}
  \item \canDien[\ldots]: số liệu đã có ở Azure Boards, dashboard chi phí hoặc biên bản họp nhưng chưa chép
        vào tài liệu; cần bổ sung trước khi trình duyệt.
  \item \textbf{Số minh hoạ}: số dùng để giải thích phương pháp (ví dụ EVA, biểu đồ burndown và đường S), không
        phải số của nhóm; chú thích của bảng và hình ghi rõ điều này.
  \item \textbf{Mã thay đổi}: CR-xx là yêu cầu thay đổi đã ghi nhận, CR-Hx là thay đổi bị hoãn hoặc từ chối;
        R-xx là mã rủi ro trong Risk Register~\cite{pmp}.
  \item \textbf{Tài liệu sống}: báo cáo được cập nhật sau mỗi sprint; số liệu Sprint 4 lấy theo repo ngày
        12/11/2026.
\end{itemize}

\tableofcontents

\chapter*{Tóm tắt}
\addcontentsline{toc}{chapter}{Tóm tắt}

Tài liệu này là \textbf{báo cáo theo dõi và kiểm soát} của dự án SumTalk: nơi tập hợp số liệu và hồ sơ sinh ra
trong lúc thực hiện, đối chiếu với đường cơ sở của kế hoạch quản lý dự án. Báo cáo gồm bốn phần, được cập nhật
sau mỗi sprint: báo cáo trạng thái từng sprint, phân tích giá trị thu được (EVA), biên bản họp và Retrospective,
và nhật ký thay đổi phạm vi (Change Log).

Sprint 4 là sprint đã ghi đầy đủ theo repo: app desktop Windows, transcript live dạng streaming (từ mới hiện sau
khoảng \textbf{1\,s}), engine STT remote và xử lý file dài đều xong; bản web đã sẵn sàng để deploy nhưng
\textbf{chưa được deploy lên AWS thật}; 88 test pytest đạt. Sự cố đáng chú ý nhất là việc push lên GitHub bị từ
chối (rủi ro R11). Change Log ghi \textbf{8 yêu cầu thay đổi} (CR-01 đến CR-08) và 3 đề xuất bị hoãn; các thay đổi
lớn nhất là thêm app desktop, live streaming, engine STT remote và bản web trên AWS. Số liệu EVA, burndown và
velocity của Sprint 1--3 đã có trên Azure Boards và dashboard nhưng chưa được chép vào; tài liệu ghi sẵn phương
pháp, công thức, cách diễn giải và một ví dụ tính đầy đủ.

\textbf{Từ khoá:} báo cáo trạng thái, sprint report, burndown, velocity, Earned Value, CPI, SPI, Retrospective,
Change Log, kiểm soát thay đổi.

\chapter*{Danh mục từ viết tắt}
\addcontentsline{toc}{chapter}{Danh mục từ viết tắt}

\begin{center}
  \setstretch{1.1}
  \begin{tabular}{@{}P{2.2cm}P{13.2cm}@{}}
    \toprule
    \thead{Viết tắt} & \thead{Ý nghĩa} \\
    \midrule
    AC     & Actual Cost -- chi phí thực tế \\
    API    & Application Programming Interface -- giao diện lập trình ứng dụng \\
    AWS    & Amazon Web Services (EC2: máy chủ ảo; S3: lưu trữ đối tượng) \\
    BAC    & Budget at Completion -- ngân sách khi hoàn thành \\
    CPI    & Cost Performance Index -- chỉ số hiệu quả chi phí \\
    CR     & Change Request -- yêu cầu thay đổi \\
    CV     & Cost Variance -- chênh lệch chi phí \\
    DoD    & Definition of Done -- tiêu chí hoàn thành \\
    EAC    & Estimate at Completion -- dự báo tổng chi phí khi hoàn thành \\
    ETC    & Estimate to Complete -- dự báo chi phí phần còn lại \\
    EV     & Earned Value -- giá trị thu được \\
    EVA    & Earned Value Analysis -- phân tích giá trị thu được \\
    GPU    & Graphics Processing Unit -- bộ xử lý đồ hoạ, dùng để chạy model \\
    LLM    & Large Language Model -- mô hình ngôn ngữ lớn \\
    PM     & Project Manager -- quản lý dự án \\
    PO     & Product Owner -- chủ sản phẩm \\
    POC    & Proof of Concept -- bản thử nghiệm khái niệm \\
    PV     & Planned Value -- giá trị kế hoạch \\
    QC     & Quality Control -- kiểm soát chất lượng \\
    SM     & Scrum Master \\
    SP     & Story Point -- điểm ước lượng độ lớn của story \\
    SPI    & Schedule Performance Index -- chỉ số hiệu quả tiến độ \\
    STT    & Speech-to-Text -- chuyển giọng nói thành văn bản \\
    SV     & Schedule Variance -- chênh lệch tiến độ \\
    TCPI   & To-Complete Performance Index -- chỉ số hiệu quả cần đạt cho phần còn lại \\
    UAT    & User Acceptance Testing -- kiểm thử chấp nhận của người dùng \\
    VAC    & Variance at Completion -- chênh lệch ngân sách khi hoàn thành \\
    \bottomrule
  \end{tabular}
\end{center}

\cleardoublepage\phantomsection
\addcontentsline{toc}{chapter}{\listfigurename}
\listoffigures

\cleardoublepage\phantomsection
\addcontentsline{toc}{chapter}{\listtablename}
\listoftables

% ============================================================ Chapter 1
\chapter{Giới thiệu}
\label{chap:intro}

\section{Bối cảnh}

Kế hoạch quản lý dự án~\cite{pmp} đặt đường cơ sở cho phạm vi (WBS), lịch trình (CPM, Gantt) và chi phí (PV theo
sprint), cùng ngưỡng hành động CPI hoặc SPI $<$ 0,9. Điều lệ dự án~\cite{pid} quy định ai được duyệt thay đổi: PO
duyệt thay đổi trong sprint, Sponsor duyệt thay đổi mốc, ngân sách vượt 10\% hoặc phạm vi dự án. Tài liệu này ghi
lại những gì \emph{thực sự} diễn ra so với đường cơ sở đó, và các quyết định kiểm soát đã đưa ra.

\section{Mục đích và người đọc của tài liệu}

Tài liệu tập hợp số liệu và hồ sơ sinh ra trong lúc thực hiện dự án: báo cáo từng sprint, phân tích giá trị thu
được, biên bản họp, Retrospective và nhật ký thay đổi phạm vi. Tài liệu được cập nhật sau mỗi sprint. Người đọc
chính là Project Manager (lập và dùng báo cáo để điều hành), Product Owner (duyệt thay đổi trong sprint, ưu tiên
lại backlog) và Sponsor (nhận bản báo cáo mỗi sprint, duyệt các thay đổi lớn).

\section{Cấu trúc tài liệu}

Phần còn lại của tài liệu gồm năm chương:
\begin{itemize}
  \item \textbf{Chương~\ref{chap:status} -- Báo cáo trạng thái theo sprint:} mẫu báo cáo, bảng tổng hợp các sprint,
        burndown và báo cáo Sprint 4.
  \item \textbf{Chương~\ref{chap:eva} -- Phân tích giá trị thu được:} định nghĩa, công thức, cách diễn giải, bảng
        theo dõi và ví dụ tính kèm đường S.
  \item \textbf{Chương~\ref{chap:minutes} -- Biên bản họp và Retrospective:} kiểm tra độ đầy đủ, mẫu biên bản bốn
        loại họp, gợi ý Retrospective Sprint 4.
  \item \textbf{Chương~\ref{chap:changelog} -- Change Log:} quy trình kiểm soát thay đổi, nhật ký thay đổi và các
        thay đổi bị hoãn.
  \item \textbf{Chương~\ref{chap:conclusion} -- Kết luận và phê duyệt.}
\end{itemize}

% ============================================================ Chapter 2
\chapter{Báo cáo trạng thái theo sprint}
\label{chap:status}

\textbf{Mỗi sprint có một báo cáo theo cùng mẫu, để Sponsor so sánh được các sprint với nhau.} Báo cáo gồm mục
tiêu, kế hoạch so với thực tế, burndown, velocity, chi phí (CPI, SPI), vấn đề, rủi ro và quyết định. Số liệu
Sprint 1--3 nhóm đã có: \canDien[chép từ báo cáo S1--S3 hiện có]. Sprint 4 được ghi theo các commit trong repo.

\section{Mẫu báo cáo sprint}

Mỗi báo cáo mở đầu bằng tiêu đề ``Sprint N -- mục tiêu sprint (ngày bắt đầu -- ngày kết thúc)'', tiếp theo là bảng
chỉ số kế hoạch so với thực tế (Bảng~\ref{tab:report-template}). Chi phí tính theo đơn giá 100.000~VND/giờ; các chỉ
số EV, CPI, SPI được định nghĩa ở Chương~\ref{chap:eva}.

\begin{table}[htbp]
  \centering
  \caption{Bảng chỉ số trong mẫu báo cáo sprint}
  \label{tab:report-template}
  \begin{tabularx}{\linewidth}{@{}LP{3.4cm}P{3.4cm}@{}}
    \toprule
    \thead{Chỉ số} & \thead{Kế hoạch} & \thead{Thực tế} \\
    \midrule
    Story point & & \\
    Số story Done / cam kết & & \\
    Giờ công & & \\
    Chi phí (giờ $\times$ 100.000 VND) & PV = & AC = \\
    EV (giá trị thu được) & & \\
    CPI = EV/AC; SPI = EV/PV & & \\
    \bottomrule
  \end{tabularx}
\end{table}

Sau bảng chỉ số, báo cáo có các mục ngắn sau:
\begin{itemize}
  \item \textbf{Burndown:} ảnh chụp từ Azure Boards, hoặc bảng điểm còn lại theo ngày (Mục~\ref{sec:burndown}).
  \item \textbf{Kết quả:} việc đã xong; việc chưa xong, kèm lý do và sprint được dời sang.
  \item \textbf{Vấn đề và rủi ro:} vấn đề phát sinh; rủi ro mới hoặc rủi ro thay đổi mức độ.
  \item \textbf{Quyết định:} các quyết định và thay đổi phạm vi (đồng thời ghi vào Change Log).
  \item \textbf{Cải tiến:} hành động cải tiến rút ra từ Retrospective.
\end{itemize}

\section{Tổng hợp các sprint}

Bảng~\ref{tab:sprints} tổng hợp bốn sprint trên một trang. Story point của Sprint 4 là số đề xuất, cộng từ năm hạng
mục: D-1 (8), S-1 (13), O-1 (8), US6.2 (5) và tài liệu (3).

\begin{table}[htbp]
  \centering
  \caption{Tổng hợp các sprint}
  \label{tab:sprints}
  \small
  \setlength{\tabcolsep}{4pt}
  \begin{tabularx}{\linewidth}{@{}Q{1.2cm}LQ{1.6cm}Q{1.2cm}Q{1.7cm}Q{0.9cm}Q{0.9cm}P{2.5cm}@{}}
    \toprule
    \thead{Sprint} & \thead{Mục tiêu} & \thead{SP cam kết} & \thead{SP xong} & \thead{Velocity} & \thead{CPI}
      & \thead{SPI} & \thead{Ghi chú} \\
    \midrule
    1 & Ghi âm, upload, transcript gốc (E1) & \canDien & & & & & \\
    2 & Transcript live, benchmark LLM & \canDien & & & & & \\
    3 & Tóm tắt \#67 và dịch \#60, test E2E & \canDien & & & & & \#60 dời vào sprint này \\
    4 & Desktop, streaming, STT remote, deploy AWS & 37 (đề xuất) & & & & & Mục~\ref{sec:sprint4} \\
    \bottomrule
  \end{tabularx}
\end{table}

\section{Burndown}
\label{sec:burndown}

Biểu đồ burndown của mỗi sprint vẽ số story point còn lại theo ngày, cạnh đường lý tưởng giảm đều về 0.
Hình~\ref{fig:burndown} là mẫu với \textbf{số minh hoạ} cho một sprint 10 ngày, 37 SP;
\canDien[thay bằng điểm còn lại theo ngày trên Azure Boards].

\begin{figure}[htbp]
  \centering
  \begin{tikzpicture}
    \begin{axis}[
        width=0.85\linewidth, height=6.4cm,
        xmin=0.6, xmax=10.4, ymin=0, ymax=40,
        xtick={1,...,10}, xticklabels={N1,N2,N3,N4,N5,N6,N7,N8,N9,N10},
        ylabel={Story point còn lại},
        grid=major, grid style={gray!25},
        legend pos=north east, legend cell align=left,
        legend style={font=\small,draw=gray!40},
        tick label style={font=\small},
        /pgf/number format/.cd, use comma, 1000 sep={.}]
      \addplot[thick, chartgray, dashed, mark=none]
        coordinates {(1,37) (2,33) (3,29) (4,26) (5,22) (6,18) (7,15) (8,11) (9,7) (10,0)};
      \addlegendentry{Lý tưởng}
      \addplot[very thick, primaryblue, mark=*]
        coordinates {(1,37) (2,37) (3,34) (4,30) (5,30) (6,24) (7,20) (8,14) (9,8) (10,0)};
      \addlegendentry{Thực tế}
    \end{axis}
  \end{tikzpicture}
  \caption{Mẫu biểu đồ burndown của một sprint (số minh hoạ)}
  \label{fig:burndown}
\end{figure}

Cách đọc: đường thực tế nằm \emph{trên} đường lý tưởng nghĩa là sprint đang chậm. Theo Risk Register~\cite{pmp},
burndown nằm trên đường lý tưởng 2 ngày liền là dấu hiệu cảnh báo của rủi ro R3 (trễ tiến độ); trong ví dụ, điều
này xảy ra từ ngày N2 và cần được nêu trong Daily.

\section{Báo cáo Sprint 4 -- desktop, streaming, cloud}
\label{sec:sprint4}

Sprint 4 được ghi theo repo ngày 12/11/2026 (Bảng~\ref{tab:sprint4}). Sáu trên bảy hạng mục đã xong; riêng deploy
AWS mới xong phần chuẩn bị.

\begin{table}[htbp]
  \centering
  \caption{Kết quả Sprint 4 theo repo (12/11/2026)}
  \label{tab:sprint4}
  \small
  \begin{tabularx}{\linewidth}{@{}P{3.2cm}LP{4.2cm}@{}}
    \toprule
    \thead{Hạng mục} & \thead{Kết quả} & \thead{Bằng chứng} \\
    \midrule
    App desktop Windows & Xong: bộ cài 204\,MB, smoke test 8/8 & Nhánh \path{feature/desktop-app} (7 commit) \\
    Live streaming & Xong: từ mới khoảng 1\,s, câu chốt trung bình 1,9\,s & Nhánh \path{feature/streaming-stt} \\
    STT remote, API \path{/v1} & Xong: máy khách không tải model, độ trễ như GPU local & Như trên \\
    File dài & Xong: file 11,6 phút xử lý trong 18--57\,s, dòng đầu sau 3--5\,s & Như trên \\
    Deploy AWS & Xong phần chuẩn bị: image, Compose, Caddy, S3, hướng dẫn; \textbf{chưa deploy lên AWS thật}
      & Nhánh \path{feature/aws-deploy} \\
    Tài liệu quản lý & 22 tài liệu \texttt{.md}, sau đó gộp thành 6 tài liệu (CR-08)
      & Nhánh \path{docs/project-management} \\
    Test & 88 test pytest đạt, ruff sạch & -- \\
    \bottomrule
  \end{tabularx}
\end{table}

\begin{itemize}
  \item \textbf{Vấn đề phát sinh:} push lên GitHub bị từ chối (403) vì tài khoản đang lưu trên máy không có quyền ghi
        repo (rủi ro R11). Các nhánh lúc đó chỉ có ở máy dev.
  \item \textbf{Quyết định:} thêm app desktop, streaming và bản web trên cloud vào phạm vi, ghi ở Change Log
        (Chương~\ref{chap:changelog}).
  \item \textbf{Việc tiếp theo:} cấp quyền repo và push; deploy thật lên AWS theo hướng dẫn; UAT với người dùng thử;
        điền số liệu giờ công Sprint 4.
\end{itemize}

% ============================================================ Chapter 3
\chapter{Phân tích giá trị thu được}
\label{chap:eva}

\textbf{EVA so sánh giá trị công việc đã làm được (EV) với chi phí thực tế (AC) và kế hoạch (PV) theo từng sprint,
rồi dự báo tổng chi phí khi hoàn thành (EAC).} Nhóm đã có CPI và SPI trên dashboard:
\canDien[chép số vào Bảng~\ref{tab:eva-track}]. Chương này ghi phương pháp, cách diễn giải và một ví dụ tính đầy đủ
với số minh hoạ.

\section{Định nghĩa và công thức}

Các đại lượng và công thức dùng ở đây theo hướng dẫn PMBOK~\cite{pmbok}. Bốn đại lượng cơ bản và cách lấy chúng
trong dự án được nêu trong Bảng~\ref{tab:eva-def}. Điểm quan trọng nhất là
EV chỉ tính các story đã đạt Definition of Done~\cite{qtd}, không tính công việc dở dang.

\begin{table}[htbp]
  \centering
  \caption{Các đại lượng cơ bản của EVA}
  \label{tab:eva-def}
  \begin{tabularx}{\linewidth}{@{}Q{1.2cm}P{4.4cm}L@{}}
    \toprule
    \thead{Ký hiệu} & \thead{Ý nghĩa} & \thead{Cách lấy trong dự án} \\
    \midrule
    BAC & Ngân sách khi hoàn thành & Tổng giờ kế hoạch $\times$ 100.000 VND (kế hoạch chi phí trong~\cite{pmp}) \\
    PV & Giá trị kế hoạch đến thời điểm đo & Giờ kế hoạch luỹ kế đến hết sprint $\times$ 100.000 VND \\
    EV & Giá trị thu được & BAC $\times$ (SP Done luỹ kế / tổng SP kế hoạch); chỉ tính story đạt DoD \\
    AC & Chi phí thực tế & Giờ thực tế luỹ kế $\times$ 100.000 VND \\
    \bottomrule
  \end{tabularx}
\end{table}

Từ bốn đại lượng trên, chênh lệch và chỉ số hiệu quả được tính theo~\eqref{eq:eva-index}, còn các dự báo khi hoàn
thành theo~\eqref{eq:eva-forecast}:
\begin{gather}
  \label{eq:eva-index}
  \mathrm{CV} = \mathrm{EV} - \mathrm{AC}, \qquad \mathrm{SV} = \mathrm{EV} - \mathrm{PV}, \qquad
  \mathrm{CPI} = \frac{\mathrm{EV}}{\mathrm{AC}}, \qquad \mathrm{SPI} = \frac{\mathrm{EV}}{\mathrm{PV}} \\
  \label{eq:eva-forecast}
  \mathrm{EAC} = \frac{\mathrm{BAC}}{\mathrm{CPI}}, \quad \mathrm{ETC} = \mathrm{EAC} - \mathrm{AC}, \quad
  \mathrm{VAC} = \mathrm{BAC} - \mathrm{EAC}, \quad
  \mathrm{TCPI} = \frac{\mathrm{BAC} - \mathrm{EV}}{\mathrm{BAC} - \mathrm{AC}}
\end{gather}

\section{Diễn giải}

Bảng~\ref{tab:eva-meaning} cho biết ý nghĩa của từng chỉ số khi lớn hơn, bằng hoặc nhỏ hơn 1, và hành động tương
ứng. Lưu ý TCPI đọc ngược chiều với CPI và SPI: TCPI lớn hơn 1 là tín hiệu xấu.

\begin{table}[htbp]
  \centering
  \caption{Diễn giải các chỉ số hiệu quả}
  \label{tab:eva-meaning}
  \small
  \begin{tabularx}{\linewidth}{@{}Q{1.2cm}P{4.3cm}P{2.6cm}L@{}}
    \toprule
    \thead{Chỉ số} & \thead{$>$ 1} & \thead{= 1} & \thead{$<$ 1} \\
    \midrule
    CPI & Tốn ít hơn kế hoạch & Đúng ngân sách & Vượt chi phí: xem lại ước lượng, cắt việc Could \\
    SPI & Nhanh hơn kế hoạch & Đúng tiến độ & Chậm: dồn người vào việc găng (lịch trình trong~\cite{pmp}) \\
    TCPI & Phải làm hiệu quả hơn trước mới kịp ngân sách & -- & Còn dư địa \\
    \bottomrule
  \end{tabularx}
\end{table}

Ngưỡng hành động của dự án giữ như trong kế hoạch chi phí~\cite{pmp}: khi CPI hoặc SPI nhỏ hơn 0,9, PM báo PO và
Sponsor.

\section{Bảng theo dõi theo sprint}

Các giá trị luỹ kế đến hết mỗi sprint được ghi vào Bảng~\ref{tab:eva-track}; các cột CV, SV, CPI, SPI, EAC tính từ
ba cột đầu theo~\eqref{eq:eva-index} và~\eqref{eq:eva-forecast}.

\begin{table}[htbp]
  \centering
  \caption{Theo dõi EVA theo sprint (luỹ kế, VND)}
  \label{tab:eva-track}
  \small
  \setlength{\tabcolsep}{4pt}
  \begin{tabularx}{\linewidth}{@{}Q{1.2cm}CCCCCQ{0.9cm}Q{0.9cm}C@{}}
    \toprule
    \thead{Hết sprint} & \thead{PV} & \thead{EV} & \thead{AC} & \thead{CV} & \thead{SV} & \thead{CPI} & \thead{SPI}
      & \thead{EAC} \\
    \midrule
    1 & \canDien & & & & & & & \\
    2 & \canDien & & & & & & & \\
    3 & \canDien & & & & & & & \\
    4 & \canDien & & & & & & & \\
    \bottomrule
  \end{tabularx}
\end{table}

\section{Ví dụ cách tính}

Để minh hoạ cách dùng công thức, giả sử \textbf{(số minh hoạ, không phải số của nhóm)} BAC = 60 triệu VND (600
giờ) và tổng 125 SP. Hết Sprint 3, kế hoạch là 450 giờ, thực tế 480 giờ, đã Done 90 SP. Các bước tính được trình bày
trong Bảng~\ref{tab:eva-example}.

\begin{table}[htbp]
  \centering
  \caption{Ví dụ tính EVA hết Sprint 3 (số minh hoạ)}
  \label{tab:eva-example}
  \begin{tabularx}{\linewidth}{@{}Q{1.4cm}P{4.6cm}L@{}}
    \toprule
    \thead{Bước} & \thead{Tính} & \thead{Kết quả} \\
    \midrule
    PV & 450 $\times$ 100.000 & 45,0 triệu \\
    EV & 60 triệu $\times$ 90/125 & 43,2 triệu \\
    AC & 480 $\times$ 100.000 & 48,0 triệu \\
    CPI & 43,2 / 48,0 & 0,90 (vượt chi phí 10\%) \\
    SPI & 43,2 / 45,0 & 0,96 (chậm 4\%) \\
    EAC & 60 / 0,90 & 66,7 triệu \\
    VAC & 60 $-$ 66,7 & $-6{,}7$ triệu (dùng gần hết 10\% dự phòng) \\
    TCPI & (60 $-$ 43,2) / (60 $-$ 48,0) & 1,40 (phần còn lại phải hiệu quả hơn 40\% mới giữ được BAC) \\
    \bottomrule
  \end{tabularx}
\end{table}

Kết luận của ví dụ: CPI chạm ngưỡng 0,9, nên PM báo PO để dời một story Should sang phiên bản sau.

Hình~\ref{fig:scurve} vẽ ba đường luỹ kế PV, EV, AC qua bốn sprint với cùng bộ số minh hoạ (giá trị hết Sprint 3
khớp với Bảng~\ref{tab:eva-example}). Khoảng cách giữa AC và EV là chênh lệch chi phí CV; khoảng cách giữa PV và EV
là chênh lệch tiến độ SV. \canDien[thay bằng số thật từ Bảng~\ref{tab:eva-track}].

\begin{figure}[htbp]
  \centering
  \begin{tikzpicture}
    \begin{axis}[
        width=0.85\linewidth, height=6.6cm,
        xmin=0.7, xmax=4.3, ymin=0, ymax=70,
        xtick={1,2,3,4}, xticklabels={Sprint 1,Sprint 2,Sprint 3,Sprint 4},
        ylabel={Triệu VND (luỹ kế)},
        grid=major, grid style={gray!25},
        legend pos=north west, legend cell align=left,
        legend style={font=\small,draw=gray!40},
        tick label style={font=\small},
        /pgf/number format/.cd, use comma, 1000 sep={.}]
      \addplot[thick, chartgray, dashed, mark=square*] coordinates {(1,15) (2,30) (3,45) (4,60)};
      \addlegendentry{PV -- kế hoạch}
      \addplot[very thick, primaryblue, mark=*] coordinates {(1,13) (2,27) (3,43.2) (4,58)};
      \addlegendentry{EV -- giá trị thu được}
      \addplot[thick, todored, mark=triangle*] coordinates {(1,14) (2,31) (3,48) (4,62)};
      \addlegendentry{AC -- chi phí thực tế}
      \draw[{Stealth[length=1.6mm]}-{Stealth[length=1.6mm]},todored] (axis cs:3.12,43.2) -- (axis cs:3.12,48)
        node[midway,right,font=\scriptsize,text=todored] {CV};
      \draw[{Stealth[length=1.6mm]}-{Stealth[length=1.6mm]},chartgray] (axis cs:2.88,43.2) -- (axis cs:2.88,45)
        node[midway,left,font=\scriptsize,text=chartgray] {SV};
    \end{axis}
  \end{tikzpicture}
  \caption{Đường S của PV, EV và AC luỹ kế (số minh hoạ)}
  \label{fig:scurve}
\end{figure}

% ============================================================ Chapter 4
\chapter{Biên bản họp và Retrospective}
\label{chap:minutes}

\textbf{Mỗi sprint cần đủ biên bản Sprint Planning, Sprint Review và Retrospective; Daily chỉ cần ghi khi có vướng
mắc.} Nhóm đã có biên bản các buổi họp: \canDien[dán liên kết hoặc chép vào Mục~\ref{sec:archive}]. Chương này đưa
ra bảng kiểm tra độ đầy đủ, mẫu biên bản cho bốn loại họp và gợi ý nội dung Retrospective Sprint 4.

\section{Kiểm tra độ đầy đủ}

Bảng~\ref{tab:minutes-check} dùng để đánh dấu biên bản đã có của từng sprint. Cột ghi chú nêu nội dung quan trọng
nhất cần có trong biên bản của sprint đó.

\begin{table}[htbp]
  \centering
  \caption{Kiểm tra độ đầy đủ của biên bản}
  \label{tab:minutes-check}
  \begin{tabularx}{\linewidth}{@{}Q{1.4cm}Q{2.2cm}Q{2.2cm}Q{2.8cm}L@{}}
    \toprule
    \thead{Sprint} & \thead{Planning} & \thead{Review} & \thead{Retrospective} & \thead{Ghi chú} \\
    \midrule
    1 & $\square$ & $\square$ & $\square$ & \\
    2 & $\square$ & $\square$ & $\square$ & Benchmark LLM; ghi rủi ro R1--R6 \\
    3 & $\square$ & $\square$ & $\square$ & Quyết định dời \#60 vào Sprint 3 \\
    4 & $\square$ & $\square$ & $\square$ & Quyết định thêm desktop, streaming, cloud \\
    \bottomrule
  \end{tabularx}
\end{table}

\section{Mẫu biên bản}

Bốn mẫu biên bản trong Bảng~\ref{tab:minutes-template} liệt kê các mục bắt buộc của từng loại họp. Sprint Planning
và Sprint Review ghi số liệu để đối chiếu với báo cáo sprint; Retrospective dùng khung Start -- Stop -- Continue và
luôn kiểm tra lại hành động của lần trước.

\begingroup\small
\setlength{\tabcolsep}{4pt}
\begin{longtable}{@{}P{4.2cm}P{10.8cm}@{}}
  \caption{Mẫu biên bản cho bốn loại họp Scrum}\label{tab:minutes-template}\\
  \toprule
  \thead{Mục} & \thead{Nội dung cần ghi} \\
  \midrule
  \endfirsthead
  \multicolumn{2}{@{}l}{\footnotesize\itshape Bảng~\thetable{} (tiếp theo)}\\
  \toprule
  \thead{Mục} & \thead{Nội dung cần ghi} \\
  \midrule
  \endhead
  \bottomrule
  \endlastfoot
  \multicolumn{2}{@{}l}{\textbf{Sprint Planning}}\\
  Thông tin chung & Sprint N Planning -- ngày -- thành phần -- chủ trì: SM \\
  Mục tiêu sprint & Một câu mô tả kết quả sprint \\
  Năng lực & Số giờ của từng người; velocity tham chiếu (SP) \\
  Story nhận vào & ID, SP, người nhận \\
  Ước lượng & Planning Poker; các story phải vote lại \\
  Rủi ro, phụ thuộc & Rủi ro mới, phụ thuộc giữa các story \\
  Quyết định & Các quyết định của buổi họp \\
  \midrule
  \multicolumn{2}{@{}l}{\textbf{Sprint Review}}\\
  Thông tin chung & Sprint N Review -- ngày -- thành phần: nhóm và Sponsor \\
  Demo & Story -- người demo -- kết quả: chấp nhận hoặc chưa \\
  Phản hồi & Phản hồi của Sponsor và PO \\
  Thay đổi backlog & Thay đổi sau review; ghi vào Change Log nếu là thay đổi phạm vi \\
  Số liệu & SP Done trên cam kết, giờ thực tế, CPI, SPI \\
  \midrule
  \multicolumn{2}{@{}l}{\textbf{Retrospective (Start -- Stop -- Continue)}}\\
  Thông tin chung & Sprint N Retrospective -- ngày \\
  Continue & Điều tốt, cần giữ \\
  Stop & Điều chưa tốt, cần bỏ \\
  Start & Điều sẽ thử mới \\
  Hành động cải tiến & Tối đa 3 hành động, mỗi hành động có người phụ trách và hạn \\
  Kiểm tra retro trước & Hành động của lần trước: đã làm hay chưa \\
  \midrule
  \multicolumn{2}{@{}l}{\textbf{Daily Scrum (chỉ ghi khi có vướng mắc)}}\\
  Vướng mắc & Ngày -- vướng mắc -- người hỗ trợ -- hạn gỡ \\
\end{longtable}
\endgroup

\section{Gợi ý nội dung Retrospective Sprint 4}

Bảng~\ref{tab:retro4} là gợi ý để nhóm thảo luận trong buổi Retrospective Sprint 4, rút ra từ diễn biến của sprint.
Hai hành động Start nhắm thẳng vào hai vấn đề của sprint: sự cố quyền push (R11) và việc bản web chưa được deploy.

% [H]: placed right after a longtable; a floating table would end up alone on its own page.
\begin{table}[H]
  \centering
  \caption{Gợi ý nội dung Retrospective Sprint 4}
  \label{tab:retro4}
  \begin{tabularx}{\linewidth}{@{}P{2.4cm}L@{}}
    \toprule
    \thead{Loại} & \thead{Nội dung} \\
    \midrule
    Continue & Chia commit theo task, mỗi nhóm việc một nhánh; đo hiệu năng bằng script thay vì ước đoán \\
    Continue & Có test tự động cho phần mới (live, STT remote, S3), chạy dưới 10 giây \\
    Stop & Để nhiều thay đổi chưa commit trong máy quá lâu (Sprint 3 chưa commit khi bắt đầu Sprint 4) \\
    Start & Kiểm tra quyền push repo ngay đầu sprint (sự cố 403 -- R11) \\
    Start & Deploy bản web sớm để QC test trên môi trường giống thật \\
    \bottomrule
  \end{tabularx}
\end{table}

\section{Lưu trữ biên bản}
\label{sec:archive}

\canDien[liên kết tới biên bản các sprint (Azure Boards Wiki, Google Drive\ldots), hoặc chép nội dung vào đây].

% ============================================================ Chapter 5
\chapter{Change Log}
\label{chap:changelog}

\textbf{Mọi thay đổi phạm vi, mốc hoặc ngân sách phải được ghi vào Change Log trước khi làm, kèm lý do, tác động và
người duyệt.} Theo phân quyền trong điều lệ dự án~\cite{pid}, thay đổi trong sprint do PO duyệt; thay đổi mốc, ngân
sách vượt 10\% hoặc phạm vi dự án do Sponsor duyệt.

\section{Quy trình kiểm soát thay đổi}

Một yêu cầu thay đổi đi qua bốn bước (Hình~\ref{fig:cr-flow}):
\begin{enumerate}
  \item Người đề xuất ghi một dòng vào nhật ký, trạng thái ``Đề xuất''.
  \item PM đánh giá tác động: story point và giờ công, chi phí, lịch (ảnh hưởng tới đường găng trong~\cite{pmp}),
        rủi ro.
  \item PO, hoặc Sponsor với thay đổi lớn, duyệt hoặc từ chối.
  \item Nếu được duyệt: cập nhật backlog, kế hoạch, ngân sách, và thông báo theo kế hoạch truyền thông~\cite{pmp}.
\end{enumerate}

\begin{figure}[htbp]
  \centering
  \begin{tikzpicture}[
      st/.style={draw=primaryblue,fill=accentblue,rounded corners=2pt,align=center,font=\footnotesize,
                 text width=2.9cm,minimum height=1.35cm,inner sep=3pt,execute at begin node=\setstretch{1.05}},
      dc/.style={draw=primaryblue,fill=white,diamond,aspect=1.6,align=center,font=\footnotesize,inner sep=1pt,
                 execute at begin node=\setstretch{1.05}},
      ar/.style={-{Stealth[length=2.2mm]},thick,primaryblue}]
    \node[st] (p) at (0,0) {\textbf{1. Đề xuất}\\ ghi một dòng vào Change Log};
    \node[st] (a) at (3.7,0) {\textbf{2. PM đánh giá}\\ SP, giờ, chi phí, lịch, rủi ro};
    \node[dc] (d) at (7.4,0) {\textbf{3. Duyệt?}\\ PO / Sponsor};
    \node[st] (u) at (12.2,0) {\textbf{4. Cập nhật}\\ backlog, kế hoạch, ngân sách; thông báo};
    \node[st,draw=chartgray,fill=gray!10] (r) at (7.4,-2.3) {Từ chối hoặc hoãn\\ (bảng CR-H)};
    \draw[ar] (p) -- (a);
    \draw[ar] (a) -- (d);
    \draw[ar] (d) -- node[above,font=\scriptsize] {duyệt} (u);
    \draw[ar,chartgray] (d) -- node[right,font=\scriptsize] {không} (r);
  \end{tikzpicture}
  \caption{Quy trình kiểm soát thay đổi}
  \label{fig:cr-flow}
\end{figure}

Quyền duyệt ở bước 3 phụ thuộc vào quy mô thay đổi: thay đổi trong một sprint do PO quyết định; thay đổi làm dời
mốc, tăng ngân sách quá 10\% hoặc mở rộng phạm vi dự án phải lên Sponsor.

\section{Nhật ký thay đổi}

Bảng~\ref{tab:changelog} ghi tám yêu cầu thay đổi theo thứ tự mới nhất trước. Người đề xuất và người duyệt của
mọi yêu cầu chưa được chép vào: \canDien[điền theo biên bản và Azure Boards].

\begingroup\small
\setlength{\tabcolsep}{4pt}
\begin{longtable}{@{}P{1.2cm}P{1.9cm}P{3.6cm}P{3.3cm}P{2.8cm}P{1.7cm}@{}}
  \caption{Change Log}\label{tab:changelog}\\
  \toprule
  \thead{Mã} & \thead{Ngày} & \thead{Thay đổi} & \thead{Lý do} & \thead{Tác động} & \thead{Trạng thái} \\
  \midrule
  \endfirsthead
  \multicolumn{6}{@{}l}{\footnotesize\itshape Bảng~\thetable{} (tiếp theo)}\\
  \toprule
  \thead{Mã} & \thead{Ngày} & \thead{Thay đổi} & \thead{Lý do} & \thead{Tác động} & \thead{Trạng thái} \\
  \midrule
  \endhead
  \bottomrule
  \endlastfoot
  CR-08 & 13/11/2026 & Gộp 22 tài liệu quản lý thành 6 tài liệu theo nhóm quy trình (khởi động, yêu cầu, kế hoạch,
    chất lượng, theo dõi, kết thúc) & Mỗi bên liên quan đọc một tài liệu thay vì nhiều file rời; dễ duyệt và xuất bản
    & Không đổi lịch, không đổi nội dung & Đã làm \\
  CR-07 & 12/11/2026 & Tài liệu quản lý dự án làm thành 21 file \texttt{.md} trong repo thay vì một tài liệu
    & Dễ chia việc, quản lý phiên bản cùng code & Không đổi lịch & Đã làm \\
  CR-06 & 12/11/2026 & Thêm bản web deploy trên AWS EC2 và S3 (Docker, Caddy, hướng dẫn) & Test và demo cho khách
    mà không cần cài app & +8 SP; chi phí AWS khoảng 2,5--3 triệu VND/tháng nếu chạy 24/7
    & Đã làm (chưa deploy thật) \\
  CR-05 & 12/11/2026 & Engine STT remote và API \path{/v1}: máy khách không cần model & Máy văn phòng không có GPU;
    giữ độ chính xác của \texttt{large-v3-turbo} & +13 SP (gộp với CR-04); thêm phụ thuộc vào server hoặc API
    & Đã làm \\
  CR-04 & 12/11/2026 & Live ``nói tới đâu hiện tới đó'' (trễ $\leq$ 3--5\,s) và xử lý file dài $\geq$ 1 giờ
    & Yêu cầu của PO sau khi dùng thử & Chunk 1\,s, LocalAgreement, chia đoạn file dài & Đã làm \\
  CR-03 & \canDien & Chuyển từ web app sang có thêm app desktop Windows & Khách muốn cài như phần mềm; thu được âm
    thanh từ Zoom/Teams bản desktop & +8 SP; thêm bộ cài, rủi ro R12 (chưa ký số) & Đã làm \\
  CR-02 & \canDien & Dời story dịch \#60 vào Sprint 3 & \canDien[lý do theo biên bản]
    & Sprint 3 nhận thêm 19 SP (\#61--\#66) & Đã làm \\
  CR-01 & \canDien & \canDien[các thay đổi trước đó trong biên bản S1--S2, ví dụ chọn model sau benchmark] & & & \\
\end{longtable}
\endgroup

Ba thay đổi mở rộng phạm vi có ghi tác động bằng story point là CR-03 (app desktop, +8 SP), CR-05 gộp CR-04
(live streaming và STT remote, +13 SP) và CR-06 (bản web trên AWS, +8 SP), tổng cộng 29 SP; CR-02 chỉ dời 19 SP đã
có trong backlog vào Sprint 3. Đây là biểu hiện cụ thể của rủi ro R5 (phình phạm vi) trong Risk Register~\cite{pmp};
các thay đổi đều được ghi lại, đúng với chiến lược ``chấp nhận có kiểm soát'' đã chọn.

\section{Thay đổi đã từ chối hoặc hoãn}

Ba đề xuất bị hoãn được giữ lại để xem xét ở thời điểm phù hợp (Bảng~\ref{tab:deferred}); chúng tương ứng với
backlog tương lai trong đặc tả yêu cầu~\cite{srs}.

% Last table of the chapter: [H] so it is not pushed onto a page of its own.
\begin{table}[H]
  \centering
  \caption{Thay đổi đã từ chối hoặc hoãn}
  \label{tab:deferred}
  \small
  \begin{tabularx}{\linewidth}{@{}P{1.4cm}P{4.6cm}LP{2.6cm}@{}}
    \toprule
    \thead{Mã} & \thead{Đề xuất} & \thead{Lý do hoãn} & \thead{Xem lại khi} \\
    \midrule
    CR-H1 & Nhận diện người nói (diarization) & Tốn thời gian, cần model riêng; không ảnh hưởng mục tiêu chính
      & Sau v1.0 \\
    CR-H2 & Tài khoản, phân quyền & Ngoài phạm vi POC; bản test dùng đăng nhập chung & Trước khi thương mại \\
    CR-H3 & Dịch trực tiếp khi đang họp (\#84), trang action items (\#79) & Ưu tiên Should/Could & Sprint sau \\
    \bottomrule
  \end{tabularx}
\end{table}

% ============================================================ Chapter 6
\chapter{Kết luận và phê duyệt}
\label{chap:conclusion}

\section{Tóm tắt}

Báo cáo đã tập hợp các hồ sơ theo dõi và kiểm soát của dự án:
\begin{itemize}
  \item \textbf{Trạng thái sprint:} mẫu báo cáo chung cho mọi sprint; Sprint 4 hoàn thành sáu trên bảy hạng mục
        (desktop, live streaming, STT remote, file dài, tài liệu, test), còn bản web chưa được deploy lên AWS thật.
  \item \textbf{EVA:} phương pháp, công thức và cách diễn giải đã sẵn sàng; ví dụ minh hoạ cho thấy cách ngưỡng CPI
        0,9 dẫn tới quyết định dời một story Should.
  \item \textbf{Biên bản:} bảng kiểm tra độ đầy đủ, mẫu biên bản bốn loại họp và gợi ý Retrospective Sprint 4.
  \item \textbf{Change Log:} 8 yêu cầu thay đổi và 3 đề xuất hoãn; ba thay đổi mở rộng phạm vi (desktop, live
        streaming và STT remote, bản web trên AWS) thêm 29 SP.
\end{itemize}

\section{Việc cần làm trước khi trình duyệt}

\begin{enumerate}
  \item Chép số liệu Sprint 1--3 từ các báo cáo hiện có vào Bảng~\ref{tab:sprints}; điền giờ công Sprint 4.
  \item Chép CPI, SPI từ dashboard vào Bảng~\ref{tab:eva-track}; thay số minh hoạ ở Hình~\ref{fig:burndown} và
        Hình~\ref{fig:scurve} bằng số thật.
  \item Đánh dấu Bảng~\ref{tab:minutes-check} và dán liên kết biên bản vào Mục~\ref{sec:archive}.
  \item Điền ngày, lý do, người đề xuất và người duyệt còn thiếu trong Change Log, nhất là CR-01 đến CR-03.
  \item Hoàn tất việc tiếp theo của Sprint 4: cấp quyền repo và push, deploy thật lên AWS, UAT với người dùng thử.
\end{enumerate}

\needspace{14\baselineskip}
\section{Phê duyệt}

Nhóm dự án đề nghị Project Manager và Product Owner xác nhận báo cáo theo dõi và kiểm soát này; bản báo cáo được gửi
Sponsor theo kế hoạch truyền thông.

\begin{table}[H]
  \centering
  \renewcommand{\arraystretch}{2.4}
  \begin{tabularx}{\linewidth}{@{}P{3.6cm}LP{3.2cm}P{2.6cm}@{}}
    \toprule
    \thead{Vai trò} & \thead{Họ tên} & \thead{Chữ ký} & \thead{Ngày} \\
    \midrule
    Project Manager & & & \\
    Product Owner & & & \\
    \bottomrule
  \end{tabularx}
\end{table}

% ============================================================ References
\begin{thebibliography}{9}
\addcontentsline{toc}{chapter}{\bibname}
  \bibitem{pid} Nhóm SumTalk, \emph{Hồ sơ khởi động dự án (Project Initiation Document)}, Tài liệu 01,
    mã SUMTALK-PID, phiên bản 2.1, 24/11/2026.
  \bibitem{srs} Nhóm SumTalk, \emph{Đặc tả yêu cầu phần mềm (Software Requirements Specification)},
    Tài liệu 02, mã SUMTALK-SRS, phiên bản 2.1, 24/11/2026.
  \bibitem{pmp} Nhóm SumTalk, \emph{Kế hoạch quản lý dự án (Project Management Plan)}, Tài liệu 03,
    mã SUMTALK-PMP, phiên bản 2.1, 24/11/2026.
  \bibitem{qtd} Nhóm SumTalk, \emph{Chất lượng và kiểm thử (Quality \& Test Document)}, Tài liệu 04,
    mã SUMTALK-QTD, phiên bản 2.1, 24/11/2026.
  \bibitem{pmbok} Project Management Institute, \emph{A Guide to the Project Management Body of Knowledge
    (PMBOK Guide)}, 6th ed., Newtown Square, PA: PMI, 2017.
\end{thebibliography}

\end{document}
